\section{CONCLUSION}\label{SC:CONCLUSION}\label{sec:conclusion}
This study presents SAD for insurance fraud screening. The original
framework combines anomaly detection, expert verification and supervised
learning in successive partitions. Its simple and multilayered structures
share the update cycle but differ in the use of entropy during initial
candidate selection. Algorithms 1--3 describe these structures explicitly.

The revised experiments compare the three architectural variants on
Carclaims and a second Insurance Claims benchmark, add a frozen-model
ablation and broaden the supervised comparisons.
On the second benchmark, Simple-B attains mean fraud F1 0.560 and AP 0.491; updating Multilayer attains F1 0.497 at review fraction 0.365. These values demonstrate an empirical workload--detection trade-off rather than a universal benefit of adding an entropy gate.

Architecture selection should therefore consider detection performance
together with the available investigation capacity.

\paragraph{Limitations and future research.}
The partitions are retrospective and random, not chronological. Stored
labels simulate expert decisions without investigation delay, errors or
monetary cost. Random audits provide verified negatives and some missed
frauds, but do not eliminate acquisition bias. The second benchmark has a
documented Kaggle access source but unverified collecting-insurer provenance.
It supports replication under a different schema and prevalence;
prospective validation at an independent insurer remains necessary.
Its smaller size also limits generalization. AP values across
different prevalences should not be interpreted as directly comparable
measures of task difficulty.

The feature screen can discard useful variables; SMOTE on encoded
categorical predictors can produce fractional values; and ranking scores
need not be calibrated probabilities. The classifier comparisons share
the Simple-B acquisition policy. Recent deep and graph-based insurance
systems requiring additional relationships or proprietary data have not
been reproduced. The adapted PSO comparator and two threshold settings do
not constitute exhaustive state-of-the-art benchmarking or optimization.

Future research should evaluate chronological claims with documented
provenance, actual review costs, matched investigation budgets,
categorical-aware resampling, calibration and prospectively fixed model
selection. The accompanying code, data, prediction records and figure
sources support reproduction of the results reported here.
